% !TeX root = ./Main.tex
% !TeX engine = xetex
% !BIB program = bibtex

\documentclass[12pt, oneside, openbib, parskip = full-]{scrartcl}
% \includeonly{ % #INFO: wenn man nur bestimmte includes nutzen will 
%     content/Anhang_KIS
%     }





\def\myStudentenname{Vorname Nachname} % #TODO: Name eintragen
\def\myTitle{Hier kommt bitte dein Title hin} % #TODO: Title eintragen
\def\myDokumententyp{Bachelorarbeit / T1000} % #TODO: Arbeitstype eintragen
%\def\myBachelorart{Bachelor of Science (B.Sc.)} % #TODO aus kommentieren wenn keine Bachelorarbeit
\def\myStudiengang{Informatik/Informationstechnik} % #TODO:Studiengang ändern
\def\myHochschule{Dualen Hochschule Baden-Württemberg Stuttgart} % #TODO: Hochschule ändern
\def\myMatrikelnummer{<Matrikelnummer>} % #TODO: Matrikelnummer ändern
\def\myKurs{<Kurses>} % #TODO: Kurs ändern
\def\myAusbildungsfirma{<FirmeName>, <Ort>}
\def\myErstgutachter{<Betreuer>} % #TODO: Erstgutachter ändern
%\def\myZweigutachter{Vorname Nachname} % #TODO aus kommentieren wenn es keinen gibt 
%\def\myTitleHeaderImage{<Image DHBW?>}
\def\myOrt{Stuttgart}
\def\Bearbeitungszeitraum{12 Wochen} % #TODO: Bearbeitungszeitraum ändern % #INFO: allgemeine Werte für die Arbeit

%%%
%%% KOMA related
%%%

\usepackage[
    paper=a4paper,
    left=25mm,
    right=25mm,
    top=25mm,
    bottom=20mm,
    includefoot,
    foot=\baselineskip,
    bindingoffset=0mm
]{geometry}

\usepackage[utf8]{inputenc} % encoding
\usepackage[T1]{fontenc} % selecting font encodings
\usepackage[ngerman]{babel} % German hyphenation
\usepackage[babel,german=quotes]{csquotes} % German quotes
\usepackage{lmodern} %MOST IMPORTANT PACKAGE!!! High resolution font
\renewcommand*\familydefault{\sfdefault} %% für Serifen freie Schrift
\usepackage{textcomp} % für in line {\textmu}

\usepackage{scrlayer-scrpage}
\clearpairofpagestyles
\setkomafont{pageheadfoot}{\sffamily\footnotesize}
\setkomafont{pagehead}{\bfseries}
\setkomafont{pagination}{}

\KOMAoptions{
    headsepline = true,
    footsepline = false,
    plainfootsepline = false
}
\renewcommand*{\sectionmarkformat}{}

\automark{section}
\ihead{\myStudentenname}
\chead{\myDokumententyp}
\ohead{\headmark}
\ofoot*{\pagemark}


\usepackage{scrdate} % Für Datum und Wochentag

%%%
%%% not KOMA related
%%%


\usepackage[onehalfspacing]{setspace}

%\usepackage{booktabs} %Never, ever use vertical rules! Never use double rules!
\usepackage{multirow}
\usepackage{comment}

%\usepackage{etoolbox} %use with care! powerful but not can create difficult errors! extends build time
%\AtBeginEnvironment{table}{\sffamily}
%\AtBeginEnvironment{tabular}{\sffamily}

%\usepackage{enumitem}
\usepackage{paralist} % for \begin{compactitem}


\usepackage{scrhack} % included because \float@listhead is not supported anymore http://www.bakoma-tex.com/doc/latex/koma-script/scrhack.pdf page 3
\usepackage{listings}
\lstset{
    basicstyle=\footnotesize, % Global Code Style
    numbers=left,
    numberstyle=\tiny,
    columns=flexible,
    breaklines=true,
    postbreak=\mbox{\textcolor{red}{$\hookrightarrow$}\space},
}
% from https://tex.stackexchange.com/questions/303465/higher-asterisks-in-lstlisting-environment
\makeatletter
\lst@CCPutMacro
    \lst@ProcessOther {"2A}{\raisebox{2pt}{*}}
    \@empty\z@\@empty
\makeatother

%\usepackage{color}
%
\definecolor{lightgray}{rgb}{0.95, 0.95, 0.95}
\definecolor{darkgray}{rgb}{0.4, 0.4, 0.4}
%\definecolor{purple}{rgb}{0.65, 0.12, 0.82}
\definecolor{editorGray}{rgb}{0.95, 0.95, 0.95}
\definecolor{editorOcher}{rgb}{1, 0.5, 0} % #FF7F00 -> rgb(239, 169, 0)
\definecolor{editorGreen}{rgb}{0, 0.5, 0} % #007C00 -> rgb(0, 124, 0)
\definecolor{orange}{rgb}{1,0.45,0.13}		
\definecolor{olive}{rgb}{0.17,0.59,0.20}
\definecolor{brown}{rgb}{0.69,0.31,0.31}
\definecolor{purple}{rgb}{0.38,0.18,0.81}
\definecolor{lightblue}{rgb}{0.1,0.57,0.7}
\definecolor{lightred}{rgb}{1,0.4,0.5}
\definecolor{grey}{rgb}{0.9,0.9,0.9}







% CSS
\lstdefinelanguage{CSS}{
  keywords={color,background-image:,margin,padding,font,weight,display,position,top,left,right,bottom,list,style,border,size,white,space,min,width, transition:, transform:, transition-property, transition-duration, transition-timing-function,background},	
  sensitive=true,
  morecomment=[l]{//},
  morecomment=[s]{/*}{*/},
  morestring=[b]',
  morestring=[b]",
  alsoletter={:},
  alsodigit={-}
}

% JavaScript
\lstdefinelanguage{JavaScript}{
  morekeywords={typeof, new, true, false, catch, function, return, null, catch, switch, var, if, in, while, do, else, case, break, function},
  morecomment=[s]{/*}{*/},
  morecomment=[l]//,
  morestring=[b]",
  morestring=[b]'
}

\lstdefinelanguage{HTML5}{
    language=html,
    sensitive=true,	
    alsoletter={<>=-},	
    morecomment=[s]{<!-}{-->},
    tag=[s],
    otherkeywords={
    % General
    >,
    % Standard tags
        <!DOCTYPE,
    </html, <html, <head, <title, </title, <style, </style, <link, </head, <meta, />,
        % body
        </body, <body,
        % Divs
        </div, <div, </div>,
        % Paragraphs
        </p, <p, </p>,
        % scripts
        </script, <script,
    % More tags...
    <canvas, /canvas>, <svg, <rect, <animateTransform, </rect>, </svg>, <video, <source, <iframe, </iframe>, </video>, <image, </image>, <header, </header, <article, </article
    },
    ndkeywords={
    % General
    =,
    % HTML attributes
    charset=, src=, id=, width=, height=, style=, type=, rel=, href=, class=, onload=,
    % SVG attributes
    fill=, attributeName=, begin=, dur=, from=, to=, poster=, controls=, x=, y=, repeatCount=, xlink:href=,
    % properties
    margin:, padding:, background-image:, border:, top:, left:, position:, width:, height:, margin-top:, margin-bottom:, font-size:, line-height:, background:, text-align:,
        % CSS3 properties
    transform:, -moz-transform:, -webkit-transform:,
    animation:, -webkit-animation:,
    transition:,  transition-duration:, transition-property:, transition-timing-function:,
    }
}

\lstdefinestyle{htmlcssjs} {%
  % General design
%  backgroundcolor=\color{editorGray},
  basicstyle={\footnotesize\ttfamily},
  frame=b,
  % line-numbers
  xleftmargin={0.75cm},
  numbers=left,
  stepnumber=1,
  firstnumber=1,
  numberfirstline=true,	
  % Code design
  identifierstyle=\color{black},
  keywordstyle=\color{blue}\bfseries,
  ndkeywordstyle=\color{editorGreen}\bfseries,
  stringstyle=\color{editorOcher}\ttfamily,
  commentstyle=\color{brown}\ttfamily,
  % Code
  language=HTML5,
  alsolanguage=JavaScript,
  alsodigit={.:;},	
  tabsize=2,
  showtabs=false,
  showspaces=false,
  showstringspaces=false,
  extendedchars=true,
  breaklines=true,
  % German umlauts
  literate=%
  {Ö}{{\"O}}1
  {Ä}{{\"A}}1
  {Ü}{{\"U}}1
  {ß}{{\ss}}1
  {ü}{{\"u}}1
  {ä}{{\"a}}1
  {ö}{{\"o}}1
}
%

%\usepackage{etex} % etex package is obsolete according to https://tex.stackexchange.com/a/186597/52414
% \#TODO: Entscheiden welchen Style
%\bibliographystyle{ieeetr} % ist ohne den Strich
%\bibliographystyle{IEEETran} % ist mit dem Strich, dafür ausführlicher
% Modifizierte Version von IEEEtran, mit Unterstützung für Online-Quellen mit URL und Datum.
% Dazu als @misc eintragen. Dann die Felder url = {...} und optional urldate = {...} einfügen.
\bibliographystyle{IEEEtran_online}

\usepackage{graphicx}
%\usepackage{svg} %no reason to use this (unless you know how to use it)
\usepackage{float}
\usepackage[abs]{overpic} % Für Unterschrift

%\usepackage[binary-units = true ]{siunitx} % option binary-units has been removed https://github.com/josephwright/siunitx/blob/main/CHANGELOG.md#removed
\usepackage{siunitx}
%\sisetup{scientific-notation = true} % Wissenschaftliche Notation
\sisetup{exponent-product = \cdot}
%\sisetup{output-product = \cdot} % Das Punkt statt X % option output-product got replaced with option product-symbol https://texdoc.org/serve/siunitx/0 page 34
\sisetup{product-symbol= \ensuremath{\cdot}}
%\sisetup{binary-units=true}
\sisetup{output-decimal-marker = {,}}
\sisetup{detect-all} 	% Einstellung das Front für SI gleich Standardfront
\sisetup{per-mode = symbol}
\sisetup{locale=DE}
\usepackage{eurosym}
\DeclareSIUnit{\sieuro}{\mbox{\euro}}

% verwende gerade kein Mathe \usepackage{amsmath, amsthm, amssymb}
% !sehr langsam  \usepackage[colorinlistoftodos,prependcaption, textsize=footnotesize]{todonotes}



%\usepackage{booktabs}
\usepackage{graphicx}
\usepackage{adjustbox}
\usepackage{subcaption}
%\usepackage{lscape}
%\usepackage{tikz}
\usepackage{pgfplots}
\pgfplotsset{compat=1.16}
\usepackage{pgfplotstable}

%% Glossar
\usepackage[acronym,nonumberlist,nopostdot]{glossaries}
\usepackage{glossary-mcols} % für Style mcolindex
\glsenablehyper % Verlinkung der Einträge im Text zu Abkürzungen
\setglossarystyle{mcolindex}
\renewcommand{\glossarypreamble}{\glsfindwidesttoplevelname[\currentglossary]} % funktioniert nicht mit mcolindex
\makenoidxglossaries


\usepackage{makecell} % Für multiline Zellen mit \makecell
%\usepackage{tabularx}
\usepackage{longtable}


\usepackage[
    hidelinks,              % Das Links nicht farbig
    bookmarks,              % Bookmarks erstellen
    bookmarksopen,          % Bookmarks bei öffnen des PDF anzeigen
    bookmarksnumbered,      % damit Nummern in Bookmarks
    bookmarksopenlevel = 1, % damit nur die erste Ebene der Bookmarks beim öffnen angezeigt werden
    pdftoolbar=		false,
    pdfmenubar=		false,
    pdftitle={		\myDokumententyp~über~\myTitle~von~\myStudentenname},
    pdfsubject={}
    pdfkeywords={}
    pdfauthor={		\myStudentenname},
    pdfstartpage={	1}
    ]{hyperref}
\addto\extrasngerman{\def\subsectionautorefname{Abschnitt}}
\addto\extrasngerman{\def\subsubsectionautorefname{Abschnitt}}
\addto\extrasngerman{\def\appendixautorefname{Anhang}} % funktioniert auch nicht
\addto\extrasngerman{\def\lstlistingautorefname{Codeausschnitt}}
\renewcommand{\lstlistingname}{Codeausschnitt}

%% Anhang
\usepackage{appendix}
\renewcommand{\appendixname}{Anhang}
\renewcommand{\appendixtocname}{Anhang}
\renewcommand{\appendixpagename}{Anhang}
\newcommand{\Appendixautorefname}{Anhang} % funktioniert nicht
 \newcommand*\appendixmore{% see the KOMA-Script documentation
   \clearpage
   %\addsec{{\appendixname}}%
   \renewcommand{\thesubsection}{{\Alph{subsection}}}%
 } % damit Appendix subsection mit Buchstaben anfangen


 \renewcommand{\k}[1]{\textit{#1}}
 \newcommand{\f}[1]{\textbf{#1}}
 \let\oldautoref\autoref
 \let\oldref\ref

 \renewcommand{\autoref}[1]{\oldautoref{#1}}
 \renewcommand{\ref}[1]{\oldref{#1}}
 \newcommand{\fullref}[1]{\autoref{#1} \nameref{#1}}
  % #INFO: Einstellungen

%\glsaddall
\newglossaryentry{control endpoint}     {name={control endpoint},       description={logische, addressierte Schnittselle in \acrshort{usb} }}
\newglossaryentry{host}                 {name={host},                   description={an example}}
\newglossaryentry{slave}                {name={slave},                  description={an example}}
\newglossaryentry{control transfer}     {name={control transfer},       description={an example}}
\newglossaryentry{bulk transfer}        {name={bulk transfer},          description={an example}}
\newglossaryentry{interrupt transfer}   {name={interrupt transfer},     description={an example}}
\newglossaryentry{isochronous transfer} {name={isochronous transfer},   description={an example}}
\newglossaryentry{pathwave}             {name={PathWave},               description={ASD}}
% bit stuffing
% descriptor
% pipe
% USB Class
% Frame
% 

\newglossaryentry{dps}{name={Doppelpulstestsystem},description={}}
\newglossaryentry{freertos}{name={FreeRTOS},description={}}
\newglossaryentry{three-state}{name={three-state},description={}}
\newglossaryentry{ieee488}{name={IEEE-488},description={}}
\newglossaryentry{ieee1394}{name={IEEE-1394 (Firewire)},description={}}
\newglossaryentry{iec60488}{name={IEC-60488},description={}}
\newglossaryentry{iec60625}{name={IEC-60625},description={}}
\newglossaryentry{tinyusb}{name={TinyUSB},description={}}


% Abkürzungen
\newacronym{usb}{USB}{\k{Universal Serial Bus}}
\newacronym{npi}{NPI}{\k{New Product Introduction}}
\newacronym{rohs}{RoHS}{\k{Restriction of Hazardous Substances Directive}}
\newacronym{awg}{AWG}{\k{American Wire Gauge}}
\newacronym{nrzi}{NRZi}{\k{Non-Return-To-Zero}}
\newacronym{gnd}{GND}{\k{Ground Potenzial}}
\newacronym{crc}{CRC}{\k{Cyclic Redundancy Check}}
\newacronym{ack}{ACK}{\k{acknowledged}}
\newacronym{dma}{DMA}{\k{Direct Memory Access}}
\newacronym{bcd}{BCD}{\k{Binary-Coded Decimal}}
\newacronym{dsug}{DSUG}{\k{Design Standard Users Group}}
\newacronym{erp}{ERP}{\k{Enterprise Resource Planning}}
\newacronym{ceq}{CEQ}{\k{Customer Experience and Quality}}
\newacronym{wince}{WinCE}{\k{Windows Embedded Compact}}
\newacronym{iso}{ISO}{\k{International Organization for Standardization}}
\newacronym{sof}{SOF}{\k{start of frame}}
\newacronym{pid}{PID}{\k{packet identifier}}
\newacronym{eop}{EOP}{\k{end of packet}}
\newacronym{nak}{NAK}{\k{not acknowledged}}
\newacronym{vid}{VID}{\k{vendor identifier}}
\newacronym{usbif}{USB-IF}{\k{\acrshort{usb} Implementers Forum}}
\newacronym{produkt_id}{PID}{\k{product identifier}}
\newacronym{iec}{IEC}{\k{International Electrotechnical Commission}}
\newacronym{ebit}{EBIT}{\k{earnings before interest and taxes}}
\newacronym{cots}{COTS}{\k{commercial-off-the-shelf}}
\newacronym{ate}{ATE}{\k{automatic test equipment}}
\newacronym{ux}{UX}{\k{user experience}}
\newacronym{kis_id}{ID}{\k{industrial design}}
\newacronym{sw}{SW}{\k{software}}
\newacronym{hw}{HW}{\k{hardware}}
\newacronym{pci}{PCI}{\k{Peripheral Component Interconnect}}
\newacronym{pcie}{PCIe}{\k{\acrshort{pci} express}}
\newacronym{pxi}{PXI}{\k{\acrshort{pci} eXtensions for Instrumentation}}
\newacronym{axie}{AXIe}{\k{Advanced Telecommunications Computing Architecture Extensions for Instrumentation and Test}}
\newacronym{odi}{ODI}{\k{Optical Data Interface}}
\newacronym{pdti}{PDTI}{\k{\acrshort{pcie} Data Transfer}}
\newacronym{hvi}{HVI}{\k{Hardware Virtual Instrument}}
\newacronym{zed}{ZED}{\k{Zero Effort Discovery}}
\newacronym{fdk}{FDK}{\k{Firmware Development Kit}}
\newacronym{fpga}{FPGA}{\k{Field Programmable Gate Array}}
\newacronym{lan}{LAN}{\k{Local Area Network}}
\newacronym{identifier}{ID}{\k{identifier}}
\newacronym{vme}{VME}{\k{Versa Module Europa bus}}
\newacronym{vxi}{VXI}{\k{\acrshort{vme} eXtensions for Instrumentation}}
\newacronym{dns}{DNS}{\k{Domain Name System}}
\newacronym{dnssd}{DNS-SD}{\k{\acrshort{dns} Service Discovery}}
\newacronym{mdns}{mDNS}{\k{Multicast \acrshort{dns}}}
\newacronym{lxi}{LXI}{\k{\acrshort{lan} eXtensions for Instrumentation}}
\newacronym{xml}{XML}{\k{Extensible Markup Language}}
\newacronym{scpi}{SCPI}{\k{Standard Commands for Programmable Instruments}}
\newacronym{hislip}{HiSLIP}{\k{High-Speed \acrshort{lan} Instrument Protocol}}
\newacronym{tcp}{TCP}{\k{Transmission Control Protocol}}
\newacronym{ddns}{DDNS}{\k{Dynamic \acrshort{dns}}}
\newacronym{mac}{MAC}{\k{Media Access Control}}
\newacronym{cms}{CMD}{\k{Code Management System}}
\newacronym{csms}{CSMS}{\k{Component and Supplier Management}}
\newacronym{tlo}{TLO}{\k{Technology Leadership Organization}}
\newacronym{ieee}{IEEE}{\k{Institute of Electrical and Electronics Engineers}}
\newacronym{gpib}{GPIB}{\k{General Purpose Interface Bus}}
\newacronym{ttl}{TTL}{\k{Transistor-Transitor Logic}}
\newacronym{srq}{SRQ}{\k{service request}}
\newacronym{eoi}{EOI}{\k{end of identify}}
\newacronym{ascii}{ASCII}{\k{American Standard Code for Information Interchange}}
\newacronym{usbtmc}{USBTMC}{\k{\acrshort{usb} Test and Measurement Class}}
\newacronym{usb488}{USBTMC-USB488}{\k{\acrshort{usbtmc} Subclass USB488}}
\newacronym{mss}{MSS}{\k{Master Summary Status}}
\newacronym{rqs}{RQS}{\k{Request Service}}
\newacronym{sre}{SRE}{\k{Service Request Enable}}
\newacronym{ese}{ESE}{\k{Event Status Enable}}
\newacronym{esr}{ESR}{\k{Event Status Register}}
\newacronym{esb}{ESB}{\k{Event Status Bit}}
\newacronym{mav}{MAV}{\k{Message avaliable}}
\newacronym{stb}{STB}{\k{Status Byte}}
\newacronym{spi}{SPI}{\k{Serial Peripheral Interface}}
\newacronym{i2c}{I\textsuperscript{2}C}{\k{Inter-Integrated Circuit}}
\newacronym{dut}{DUT}{\k{device unter test}}
\newacronym{eeprom}{EEPROM}{\k{electrically erasable programmable read-only memory}}
\newacronym{mosfet}{MOSFET}{\k{metal–oxide–semiconductor field-effect transistor}}
\newacronym{hal}{HAL}{\k{Hardware Abstraction Layer}}
\newacronym{ll}{LL}{\k{Low-Layer}}
\newacronym{http}{HTTP}{\k{Hypertext Transfer Protocol}}
\newacronym{mit}{MIT}{Massachusetts Institute for Technology}
\newacronym{ansi}{ANSI}{American National Standards Institute}

\newacronym{visa}{VISA}{\k{Virtual Instrument Software Architecture}}

\newacronym{devOut}{DEV\_DEP\_MSG\_OUT}{\k{device dependent command message}}
\newacronym{devIn}{DEV\_DEP\_MSG\_IN}{\k{device dependent response message}}
\newacronym{vendorOut}{VENDOR\_SPECIFIC\_OUT}{\k{vendor specific command message}}
\newacronym{vendorIn}{VENDOR\_SPECIFIC\_IN}{\k{vendor specific response message}}

\newacronym{msgid}{MsgID}{\k{message identifier}}
\newacronym{mep}{MEP}{\k{Message Exchange Protocol}}

\newacronym{com}{COM}{\k{Component Object Model}}

\newacronym{usbmsc}{USB-MSC}{USB \k{mass storage device class}}
\newacronym{usbcdc}{USB-CDC}{USB \k{communication device class}}
\newacronym{usbhid}{USB-HID}{USB \k{human interface device class}}
\newacronym{usbmidi}{USB-MIDI}{USB \k{musical instrument digital interface device class}}
\newacronym{usbmtp}{USB-MTP}{USB \k{media transfer protocol device class}}
\newacronym{fat}{FAT}{File Allocation Table}
\newacronym{scsi}{SCSI}{Small Computer System Interface}
\newacronym{sram}{SRAM}{Static Random-Access Memory}
\newacronym{hdd}{HDD}{Hard Disk Drive}
 % #INFO: laden aller Glossary und Abkürzungen

\begin{document}
%%%
%%% Vorbau
%%%
\pagenumbering{Roman}
\begin{titlepage}
	\thispagestyle{empty}
    \begin{center}
        \begin{minipage}{6in}
            \centering
            \raisebox{-0.5\height}{\includegraphics[width=0.35\textwidth]{content/images/DHBW-Logo.pdf}}
            \hfill
            \raisebox{-0.5\height}{\includegraphics[width=0.4\textwidth]{content/images/KY_Logo_New.pdf}}
        \end{minipage}
        \vspace{1cm}
        \\
        \rule{\textwidth}{1.6pt}\vspace*{-\baselineskip}\vspace*{2pt}
        \rule{\textwidth}{0.4pt}
        \\[0.2\baselineskip]
        {\LARGE \myTitle}
        \rule{\textwidth}{0.4pt}\vspace*{-\baselineskip}\vspace{3.2pt}
        \rule{\textwidth}{1.6pt}
        \\
        \vspace*{12mm}	{\large \myDokumententyp}\\
        \ifx\myBachelorart\undefinde
            \vspace*{27mm}	des Studiengangs \myStudiengang\\
        \else
            \vspace*{12mm}	für die Prüfung zum\\
            \vspace*{3mm} 	{\large \myBachelorart}\\
            \vspace*{12mm}	Studienganges \myStudiengang\\
        \fi
        \vspace*{3mm} 	an der \myHochschule\\
        \vspace*{12mm}	von\\
        \vspace*{3mm} 	{\large \myStudentenname}\\
        \vspace*{12mm}	{\todaysname}, der {\today} \\
      \end{center}

      \vfill
      \begin{spacing}{1.2}
        \centering
        \begin{tabbing}
            \hspace{6.5cm}                   \= \kill\\
            \textbf{Matrikelnummer, Kurs}  \>  \myMatrikelnummer, \myKurs\\
            \textbf{Ausbildungsfirma}      \>  \myAusbildungsfirma\\
            \textbf{Betreuer}         \>  \myErstgutachter\\
            \ifx\myZweigutachter\undefined
            \else
            \textbf{Zweitgutachter}        \>  \myZweigutachter\\
            \fi
            \textbf{Bearbeitungszeitraum}  \>  \Bearbeitungszeitraum\\
        \end{tabbing}
      \end{spacing}
\end{titlepage}

\section*{Selbständigkeitserklärung}
\vspace{1cm}
\begin{tabbing}
    \hspace{.175\textwidth} \= \hspace{.825\textwidth}\=\kill\\
    Name:		            \> \myStudentenname\\
    Matrikelnummer:		    \> \myMatrikelnummer\\
    Studiengang:		    \> \myStudiengang\\
    Kurs:				    \> \myKurs\\
    Titel:	                \> \myTitle             \>\\
\end{tabbing}
\vspace{1cm}
Ich versichere hiermit, dass ich die vorliegende Arbeit selbstständig verfasst und
keine anderen als die angegebenen Quellen und Hilfsmittel benutzt habe.\\
Falls sowohl eine gedruckte als auch elektronische Fassung abgegeben wurde,
versichere ich zudem, dass die eingereichte elektronische Fassung mit der gedruckten
Fassung übereinstimmt.\\
\vspace{2cm}\\
\begin{minipage}{0.99\textwidth}
	\centering
	\begin{minipage}[t]{0.5\textwidth}
	\hspace{2cm}
		\begin{tabular}{c}
			\myOrt, {\today}\\
			\\
			\\
		\end{tabular}
	\end{minipage}
	\hfill
	\begin{minipage}[t]{0.4\textwidth}
		\begin{tabular}{c}
			{\myStudentenname}\\
			\\
			\\
		\end{tabular}
	\end{minipage}
\end{minipage}

\section*{Abstract}
\#TODO: Abstract
\section*{Kurzfassung}
\#TODO: Kurzfassung

% --- Inhaltsverzeichnis ---
\pdfbookmark[section]{\contentsname}{toc}
\tableofcontents

\clearpage
% --- Abbildungsverzeichnis --
\phantomsection
\addcontentsline{toc}{subsection}{Abbildungsverzeichnis}
\listoffigures

% --- Tabellenverzeichnis ---
\phantomsection
\addcontentsline{toc}{subsection}{Tabellenverzeichnis}
\listoftables

% --- Listing Verzeichnis --- % #INFO: Verwende ich gerade nicht
\phantomsection
\addcontentsline{toc}{subsection}{Codeausschnittsverzeichnis}
\lstlistoflistings

% --- Glossar ---
\iffalse % #INFO: auf \iftrue setzen falls man ein Glossar will
\phantomsection
\addcontentsline{toc}{subsection}{Glossar}
\printnoidxglossary[sort=word,title=Glossar] % main glossary
\fi

% --- Abkürzungen ---
\phantomsection
\addcontentsline{toc}{subsection}{Abkürzungen}
\printnoidxglossary[type=acronym,title=Abkürzungen] % Abkürzungen


\clearpage

%%%
%%% Inhalt
%%%
\pagenumbering{arabic}
% #INFO: für Inhalt normalerweise \input. \include with care, kann Verweise und so schwieriger machen
% ich nutze trotzdem include damit ich mit \includeonly{} zum test bauen kleiner PDFs erstellen lassen kann
%\input{content/<Name des Kapitels>}
%\include{content/<filename>}

\section{Vorlage}
\subsection[USB]{\acrshort{usb}}
\acrfull{usb} ist ein Schnittstelle zur Anbindung von \acrshort{usb}-Geräten.


\subsubsection{Zitieren} \label{sssec:zitieren}

Zitieren \cite[s. 31]{usb_developer_guide} angeschlossen werden, die sich auf Stern-Sternnetzten, mit zu 7 Ebenen \cite[4.1.1]{usb_developer_guide}, aufteilen können. Der Steckplatz von einem Hub zum Host, als upstream port bezeichnet, Buchsen für Peripheriegeräten oder weiter Hubs als downstream ports \cite{AboutBlitzerDe}.
Die Struktur einer logischen Verbindung von Host zu einem Slave ist in \autoref{fig:usb_pipe_modell} 
\begin{figure}[h]
    \centering
    \includegraphics[width=0.75\textwidth,height=6cm,keepaspectratio=true]{content/images/Cypress_ubs_figure_1.PNG}
    \caption{\acrshort{usb} pipe Modell \cite[Figure~1]{usb_developer_guide}}
    \label{fig:usb_pipe_modell}
\end{figure}
\par
Ein endpoints kann entweder die erwähnten Protokollfunktionen bereitstellen, Daten vom Host zu empfangen, dann ist dies ein IN endpoint, oder Daten für den Host bereitzustellen dann ist dies ein OUT endpoint. Diese Terminologie folgt der in der USB Spezifikation verwendete Betrachtung des Bus aus Sicht des Hosts.
Neben der Richtung eines endpoints, IN oder OUT, werden vier Übertragungsarten \cite[5.4]{usb_developer_guide} unterschieden. Die Übertragungsart wird durch die Gerätekonfiguration für jeden endpoint festleget und wird in \autoref{sssec:endpoints} weiter ausgeführt. Die Funktionen der Übertragungsarten sind \cite[4.7]{usb_developer_guide}:

\begin{compactitem}
    \item \textbf{control transfer} Für das Übertragen von Konfigurations- und Steuerungsbefehlen über die control pipe.
    \item \textbf{interrupt transfer} Für die Übertragung kleiner Datenmengen mit einer möglichst kleinen Latenz über eine data pipe .
    \item \textbf{bulk transfer} Für die Übertragung grö{\ss}er Datenmengen ohne zugesicherten Geschwindigkeit oder Latenz über eine data pipe ohne Datenverluste.
    \item \textbf{isochronous transfer}  Für Übertragungen mit zugesicherten Übertragungskapazität über eine data pipe ohne Fehlerbehebung.
\end{compactitem}
Welche Übertragungsarten ein Gerät unterstützt, und damit über welche endpoints es verfügt kann frei gewählt werden oder kann einer %\acrshort{ubs} Klassen Spezifikation folgen. Eine solche Spezifikation definiert neben dem Aufbau der endpoints weiter einheitliche Funktionen damit Geräte, einer gleichen Klasse, problemlos ausgetauscht werden können.

Softwareseitig wird die Übertragung in drei Ebenen unterteilt, diese sind in \autoref{fig:usb_schnittstellen_abstraktion} abgebildet.
\begin{figure}[h]
    \centering
    \includegraphics[width=.75\textwidth,keepaspectratio=true]{content/images/usb_spec_figure_5_2.PNG}
    \caption{\acrshort{usb} Schnittstellen Abstraktion \cite[Figure~5-2]{usb_developer_guide}}
    \label{fig:usb_schnittstellen_abstraktion}
\end{figure}



\subsubsection{Endpoints} \label{sssec:endpoints}
Nach  \acrshort{usb} Spezifikation ist ein endpoint eine eindeutiger adressierbarer Schnittstelle eines \acrshort{usb} Gerätes der Daten senden oder empfangen werden. Ausgenommen ist der control endpoint, dessen Verwendung wird in \autoref{sssec:Konfigurationsvorgang} ausführlicher beschrieben.

Die \acrshort{usb} Spezifikation definiert vier Arten von endpoints mit der maximalen Paketgrö{\ss}e und den unterstützen Übertragungsgeschwindigkeiten. Die Art eines endpoints definiert gleichzeitig die die \k{transfer}-Art aus \fullref{sssec:usb_ueberblick} und sind \cite[Kapitel~8]{usb_developer_guide}:
\begin{compactitem}
    \item \textbf{control endpoint} für control Transfers und muss in jedem Gerät vorhanden sein, um Geräteinformation zu übertragen. Die Übertragung sind fehlerbehebend und können, bei High-Speed, bis zu 10\% der Buslast einnehmen.
    \item \textbf{control endpoint} für control Transfers und muss in jedem Gerät vorhanden sein, um Geräteinformation zu übertragen. Die Übertragung sind fehlerbehebend und können, bei High-Speed, bis zu 10\% der Buslast einnehmen.
\end{compactitem}

Im  device descriptor werden allgemeine Informationen für ein Gerät bereitgestellt, die Datenfelder sind in \autoref{tab:ubs_device_descriptor} aufgetragen.
\begin{table}[h]
    \centering
    \begin{tabular}{|l|l|}
        \hline
        \textbf{Feld}      & \textbf{Bedeutung}                                 \\
        \hline
        bLength            & Länge des descriptor               \\
        \hline
        bDescriptorType    & descriptor Type, hier DEVICE (0x01)                 \\
        \hline
        bcdUSB             & \acrshort{usb} Spezifikationsversion als \acrshort{bcd}                   \\
        \hline
        bDeviceClass       & device class                                       \\
        \hline
        bDeviceSubClass    & device subclass                                    \\
        \hline
        bDeviceProtocol    & Device Protocol                                    \\
        \hline
        bMaxPacketSize     & Maximale Packet Größe für endpoint 0               \\
        \hline
        idVendor           & \acrfull{vid}, vergeben durch \acrshort{usbif})        \\
        \hline
        idProduct          & \acrfull{produkt_id}, vergeben durch den Hersteller   \\
        \hline
        bcdDevice          & Geräteversion als \acrshort{bcd}                              \\
        \hline
        iManufacturer      & Position der Herstellerbezeichnung in Packet       \\
        \hline
        iProduct           & Position der Gerätebezeichnung in Packet           \\
        \hline
        iSerialNumber      & Position der Seriennummer in Packet                \\
        \hline
        bNumConfigurations & Anzahl der unterschiedlichen Gerätekonfiguration   \\
        \hline
    \end{tabular}
    \caption{device descriptor}
    \label{tab:usb_device_descriptor} % #TODO Offset und Size entfernen
\end{table}



    %\include{content/Schluss}

%%%
%%% Nachbau
%%%
\addsec{Literatur}
\renewcommand{\section}[2]{}
\bibliography{citation/quellen}
\clearpage

% Anhänge
\appendix
\addappheadtotoc
\appendixpage
\begin{appendices}
%\include{content/<Name des Anhangs>}
\newcommand*{\tabletype}{longtable}

\subsection[KIS Requirements]{Requirements }\label{app:anhang_kis_requirements}
\subsubsection{Instrument Hardware ID Requirements}
\begin{\tabletype}[H]{|p{.12\linewidth}|p{.83\linewidth}|}
\hline
\textbf{Nr.} & \textbf{Requirement} \\ \hline 
\endhead 
4.1-2.3.1 & Existing  devices MAY continue to use the existing JTAG or JEDEC manufacturer ID \\ \hline 
4.1-2.3.2 & New  JEDEC devices (released after August 1, 2014) MUST use the  JEDEC manufacturer ID \\ \hline 
\caption {Instrument Hardware ID Requirements}
\label{tab:InstrumentHardwareIDRequirementsList}
\end{\tabletype}


\end{appendices}
\end{document}
